\section{The group quotient in cote 123: central stabilizers and an exact type-A rank formula}

The rank question in the letter, cote 123 physical page 23, footer 22,
mentions both the Lie-algebra quotient and $G\to T/W$ for the original
complex adjoint group.  The simply connected qualification occurs later,
for rational singularities.  These scopes matter: a regular element of
an adjoint group can have quotient differential of rank zero.  The
following complete deductions give a finite formula in every semisimple
type and a closed matrix formula in type A.

\subsection{BR-12305. The differential of a finite linear quotient}
Let a finite group $Z$ act linearly on a finite-dimensional complex
vector space $V$, let $\delta:V\to V/Z$ be the affine quotient, and
let $H=Z_x$ be the stabilizer of $x\in V$.  Identify $T_xV$ with $V$.
Then
\[
 \operatorname{im}(d\delta_x)^*=(V^*)^H.
 \tag{GR.1}
\]
Consequently, for any linear map $L:U\to V$, the rank of
$d\delta_x\circ L$ is the rank of the restrictions to $L(U)$ of
the $H$-invariant linear forms on $V$.

\paragraph{Complete proof.}
Every invariant polynomial has an $H$-invariant differential at $x$,
since $H$ fixes $x$.  This proves one inclusion in (GR.1).
Conversely, let $\lambda\in(V^*)^H$.  The maximal ideals of the
finitely many distinct points of $Zx$ are pairwise comaximal, as are
their squares.  The Chinese remainder theorem supplies a polynomial
$e$ with $e\equiv1\pmod{\mathfrak m_x^2}$ and
$e\equiv0\pmod{\mathfrak m_y^2}$ for every $y\in Zx$, $y\ne x$.
Put
\[
 f(v)=\frac1{|H|}\sum_{z\in Z}
       \lambda(z^{-1}v-x)e(z^{-1}v).
\]
Reindexing the sum shows that $f$ is $Z$-invariant.  At $v=x$,
the terms with $z\notin H$ have zero first jet.  A term with
$z\in H$ has differential $\lambda\circ z^{-1}=\lambda$;
its factor $\lambda(z^{-1}x-x)$ is zero, so differentiating $e$
adds nothing.  Therefore $df_x=\lambda$.  This proves equality.
The image of the dual of a composite is the restriction of this
image along $L$, so its dimension gives the final assertion.
\(\square\)

No smoothness of $V/Z$ at $\delta(x)$ was assumed.  Its Zariski
tangent space can have dimension larger than $\dim V$; (GR.1)
still computes the actual differential rank.

\subsection{BR-12306. A finite character formula for every central quotient}
Let $\widetilde G$ be a connected simply connected complex
semisimple group of rank $r$, let $Z$ be a subgroup of its center,
and put $G=\widetilde G/Z$.  For the fundamental representations
$\rho_i$ of $\widetilde G$, write
$\chi_i=\operatorname{Tr}\rho_i$ and let $\xi_i:Z\to\mathbb C^*$
be their central characters.  Steinberg (1965), Theorem 6.1(b),
p.~62, and the quotient discussion (6.9)--(6.10), pp.~63--64,
give the polynomial coordinates
\[
 \widetilde\chi=(\chi_1,\ldots,\chi_r):
 \widetilde G\longrightarrow\mathbb A^r
       =\widetilde T/W.
\]
For $g\in G$, choose a lift $\widetilde g$, put
$x_i=\chi_i(\widetilde g)$, and define
\[
 H_x=\{z\in Z:\xi_i(z)x_i=x_i\text{ for every }i\},
 \qquad J_x=\{i:\xi_i|_{H_x}=1\}.
 \tag{GR.2}
\]
If $\chi_G:G\to T/W$ is the group quotient, then
\[
 \operatorname{rank}d\chi_G|_g
 =\dim_{\mathbb C}\operatorname{span}
 \left\{Y\longmapsto
 \operatorname{Tr}\bigl(\rho_i(\widetilde g)d\rho_i(Y)\bigr)
 :i\in J_x\right\}\subseteq\mathfrak g^*.
 \tag{GR.3}
\]
This is an explicit finite formula: evaluate the fundamental
characters, determine the finite stabilizer (GR.2), and compute
the rank of the displayed $|J_x|$ linear forms on a Lie-algebra
basis.  It applies to arbitrary elements, including those with
nontrivial unipotent part.

\paragraph{Complete proof.}
The central isogeny $\widetilde G\to G$ is \'etale in characteristic
zero, and its differential is an isomorphism.  Its action on the
fundamental-character coordinates is diagonal with weights $\xi_i$,
because a central element acts as the scalar $\xi_i(z)$ in an
irreducible representation.  Taking invariants in either order gives
\[
 \mathbb C[G]^G
 =\bigl(\mathbb C[\widetilde G]^{\widetilde G}\bigr)^Z
 =\mathbb C[x_1,\ldots,x_r]^Z.
\]
Thus the lifted quotient map is the composite of
$\widetilde\chi$ with $\delta:\mathbb A^r\to\mathbb A^r/Z$.
The $H_x$-invariant linear forms are precisely the span of $dx_i$
with $i\in J_x$.  BR-12305 gives the rank as that of these
coordinate differentials composed with $d\widetilde\chi$.

Identify $T_{\widetilde g}\widetilde G$ with $\mathfrak g$ by
$Y\mapsto\widetilde gY$.  Differentiating
$\chi_i(\widetilde g\exp(tY))$ at zero gives exactly the trace
in (GR.3).  Replacing the lift by $z\widetilde g$ leaves $H_x$
and $J_x$ unchanged and multiplies each displayed form by the
nonzero scalar $\xi_i(z)$.  The rank is therefore independent
of the lift.  This completes the proof. \(\square\)

For a regular lift, Steinberg's Theorem 8.1, p.~70, makes all
$d\chi_i$ independent.  In that case (GR.3) simplifies to
$\operatorname{rank}d\chi_G|_g=|J_x|$.  In particular the
simply connected case has full rank $r$ at every regular element,
whereas the adjoint quotient can have a smaller rank.

\subsection{BR-12307. Closed formula for all type-A central quotients}
Let $c$ divide $n$, let $G=\operatorname{SL}_n/\mu_c$, and let
$A\in\operatorname{SL}_n(\mathbb C)$ lift $g\in G$.
Write its characteristic polynomial as
\[
 \det(tI-A)=\sum_{i=0}^{n}(-1)^i c_i(A)t^{n-i},
 \qquad c_0=c_n=1.
\]
Define
\[
 d=\gcd\bigl(c,\{i:1\leq i<n,\ c_i(A)\ne0\}\bigr),
 \tag{GR.4}
\]
with the convention that the gcd of the empty index set and $c$
is $c$.  If $m_{A^d}$ is the minimal polynomial of $A^d$, then
\[
 \boxed{\operatorname{rank}d\chi_G|_g
        =\deg m_{A^d}-1.}
 \tag{GR.5}
\]
The formula covers $c=1$ (the simply connected group) and $c=n$
(the adjoint group $\operatorname{PGL}_n$).

\paragraph{Complete proof.}
For $\operatorname{SL}_n$, the fundamental characters are
$c_i=\operatorname{Tr}(\bigwedge^i A)$, $1\leq i<n$.
The scalar $\zeta I\in\mu_c$ acts with central character
$\zeta^i$.  Hence (GR.2) gives $H_x=\mu_d$ and
\[
 J_x=\{d,2d,\ldots,n-d\}.
\]
All nonzero $c_i$ have index divisible by $d$, and $d\mid n$,
so $\det(tI-A)=f(t^d)$ for a monic polynomial $f$ of degree
$h=n/d$.  The Cayley--Hamilton identity gives $f(A^d)=0$.
In particular $\deg m_{A^d}\leq h$.

To compute the differentials without an eigenvalue-genericity
assumption, differentiate the identity
$\det(I+tA)=\sum_i c_i t^i$ in the tangent direction $AY$.
The determinant derivative and the formal geometric series give
\[
 dc_i|_A(AY)=\sum_{j=1}^{i}(-1)^{j-1}c_{i-j}(A)
                                  \operatorname{Tr}(A^jY).
 \tag{GR.6}
\]
For $i=md$, a nonzero coefficient on the right forces $d\mid j$.
Writing $B=A^d$, the forms $dc_{md}$, $1\leq m<h$, are
therefore related to
$Y\mapsto\operatorname{Tr}(B^jY)$, $1\leq j<h$, by a
triangular matrix whose diagonal entries are
$(-1)^{md-1}$, and are nonzero.

The annihilator of $\mathfrak{sl}_n$ in
$\mathfrak{gl}_n$ under the trace pairing is the scalar line.
The dimension of these forms restricted to $\mathfrak{sl}_n$
is consequently
\[
 \dim\operatorname{span}\{I,B,\ldots,B^{h-1}\}-1
 =\deg m_B-1.
\]
The last equality uses $\deg m_B\leq h$ and the defining
minimality of the minimal polynomial.  Equation (GR.3) now
proves (GR.5).  No Jordan block was discarded and no
semisimplicity of $A$ was used. \(\square\)

Equivalently, $\deg m_{A^d}$ is the sum, over distinct nonzero
eigenvalues $\eta$ of $A^d$, of the largest Jordan-block size
among eigenvalues $\lambda$ of $A$ with $\lambda^d=\eta$.
Indeed the polynomial $t^d$ has nonzero derivative at each
nonzero $\lambda$, so it preserves the sizes of the local
Jordan blocks, while it can merge their eigenvalue labels.
This gives a direct formula from the Jordan data of $A$.

\subsection{BR-12308. Regular adjoint elements with rank-zero quotient differential}
For every $n\geq2$, the group $\operatorname{PGL}_n(\mathbb C)$
has a regular semisimple element $g$ with
$\operatorname{rank}d\chi_G|_g=0$.
In the simply connected group $\operatorname{SL}_n$, the
corresponding regular lift has rank $n-1$.

\paragraph{Complete counterexample.}
Fix a primitive $n$th root of unity $\zeta$, set
$Q=\operatorname{diag}(1,\zeta,\ldots,\zeta^{n-1})$, and
take $g=[Q]$.  Its eigenvalues are distinct, so its connected
centralizer is a maximal torus of dimension $n-1$.  It is
therefore regular semisimple.  Choose a scalar $a\ne0$ with
$A=aQ\in\operatorname{SL}_n$.  Its characteristic polynomial
is $t^n-a^n$, so $c_i(A)=0$ for every $1\leq i<n$.
For $G=\operatorname{PGL}_n$ one has $c=n$, hence $d=n$ in
(GR.4), and $A^n=a^n I$ has minimal polynomial of degree one.
Equation (GR.5) gives rank zero.  For the simply connected
group one instead has $c=d=1$, and the distinct eigenvalues
give $\deg m_A=n$, hence rank $n-1$.
\(\square\)

This example concerns the finite central quotient of the
\emph{group} invariant space.  It does not contradict Kostant's
regular-point criterion for the Lie-algebra quotient or Steinberg's
simply connected group theorem.  The original rank questions
must therefore keep their two quotient maps separate.

\paragraph{Primary input.}
Steinberg, R. (1965). Regular elements of semi-simple algebraic
groups. \emph{Publications Math\'ematiques de l'IH\'ES, 25},
49--80. \url{https://doi.org/10.1007/BF02684397}.
Theorem 6.1(b), p.~62; (6.9)--(6.10), pp.~63--64;
Theorem 8.1, p.~70. The finite-quotient differential lemma and
the formulas (GR.3)--(GR.6) are completely proved above.
